\documentclass{article}

\usepackage[utf8]{inputenc}
\usepackage[T1]{fontenc}
\usepackage{helvet}
\usepackage{geometry}
\usepackage{xcolor}
\usepackage{eso-pic}
\usepackage{fancyhdr}
\usepackage{parskip}
\usepackage{microtype}
\usepackage[english]{babel}
\usepackage{url}
\usepackage{amsmath}
\usepackage{amssymb}
\usepackage{amsthm}

\definecolor{nato}{HTML}{293757}
\definecolor{footergray}{gray}{0.65}

\geometry{a4paper,left=3cm,right=3cm,top=3cm,bottom=3cm}
\renewcommand{\familydefault}{\sfdefault}
\setcounter{secnumdepth}{-2}

\let\oldsection\section
\renewcommand\section[1]{
  \vspace{30pt}{\color{nato}\Huge\bfseries #1\par}
  \vspace{12pt} \addcontentsline{toc}{section}{#1}}

\let\oldsubsection\subsection
\renewcommand\subsection[1]{
  \vspace{20pt}{\color{nato}\LARGE\bfseries #1\par}
  \vspace{8pt}\addcontentsline{toc}{subsection}{#1}}

\fancyhf{}
\fancyhead[R]{\small\color{footergray} March 03, 2026}
\fancyfoot[L]{\small\color{footergray} PureScript vs Haskell. Pure Functional Concurrent Language for Erlang/OTP. Interoperability and Elixir Reimplementation of Hindley-Milner Language for Production}
\fancyfoot[R]{\small\color{footergray}\thepage}

\newcommand\HUGE[1]{\fontsize{40}{40}\selectfont #1}

\renewcommand{\headrulewidth}{0pt}
\renewcommand{\footrulewidth}{0.4pt}
\renewcommand{\footrule}{\color{footergray}\hrule width \textwidth height 0.4pt}

\begin{document}

\pagestyle{fancy}

\begin{titlepage}
  \AddToShipoutPictureBG*{\AtPageLowerLeft{\color{nato}\rule{\paperwidth}{\paperheight}}}
  \color{white}\thispagestyle{empty}\vspace*{4cm}\centering
  {\LARGE \textbf{PureScript vs Haskell} \par} \vspace{2cm}
  {\HUGE  \textbf{Pure Functional Concurrent Language for Erlang/OTP} \par} \vspace{0.8cm}
  {\LARGE \textbf{Interoperability and Elixir Reimplementation of Hindley-Milner Language for Production} \par} \vspace{2.5cm}
  {\large Technical Article --- The Phi Project \par} \vfill
  {\small Version 1.0. Published March 03, 2026 \par}
  {\small \url{https://github.com/synrc/phi} \par}
\end{titlepage}

\newpage

\begin{center}
  {\Huge\color{nato}\bfseries Contents\par}
\end{center}
\vspace{40pt}
\tableofcontents

\begin{abstract}
  PureScript and Haskell offer powerful Hindley-Milner type systems and pure functional paradigms that can replace classical imperative languages in production.
  However, neither natively targets the Erlang VM (BEAM), which excels in massive concurrency, fault tolerance, distribution, and soft real-time applications (telecom, IoT, 5G edge, etc.).
  This article compares PureScript and Haskell with a focus on their type-system internals and production trade-offs, but reframes the discussion around the ultimate goal: achieving full BEAM interoperability through a PureScript-like language fully reimplemented in Elixir (as pursued by the Phi project at https://github.com/synrc/phi).
  Phi combines HM-style static typing and functional expressiveness with native BEAM primitives, avoiding JavaScript or native-code limitations while inheriting Elixir/Erlang's battle-tested concurrency model.
\end{abstract}

\newpage

\section{Production Properties of BEAM}

\subsection{Introduction}

PureScript and Haskell both build on Hindley-Milner (HM) inference for strong static typing, parametric polymorphism, and principal types. They excel at compile-time safety compared to classical languages.

Yet for systems requiring actor-based concurrency, supervision trees, hot code swapping, and seamless distribution across nodes — hallmarks of Erlang/OTP and BEAM — neither is optimal out of the box:

\begin{itemize}
\item PureScript targets JavaScript runtimes (Node/browser), lacking native BEAM integration which was fixed lately with Hamler but not finishd. We took our place in game.
\item Haskell offers GHCJS, Fay or external FFI bridges, but loses BEAM's lightweight processes and OTP patterns,
      Cloud Haskell library tries to mimic Erlang runtime but fails as Scala Akka in terms of performance, simplicity, and feature completeness.
\end{itemize}

A reimplementation in Elixir (like Phi) preserves PureScript-like syntax and HM typing while compiling directly to BEAM bytecode, unlocking Erlang's concurrency primitives natively.

PureScript remains strict and extension-free; Haskell lazy and extension-heavy. Phi inherits PureScript's simplicity while adding BEAM-native distribution.

\subsection{Production Comparison Table}

\begin{tabular}{|l|p{4.8cm}|p{4.8cm}|}
\hline
Aspect & Haskell & PureScript \\
\hline
Maturity \& Ecosystem & Very mature; thousands of packages; used in finance, AI, large backends & Younger; smaller ecosystem; excellent JS interop; strong in frontend (Halogen, React) \\
Performance & Native binaries; high throughput; laziness optimizes certain patterns & JS runtime; strict evaluation; predictable but slower than native \\
Interop \& Deployment & C FFI; GHCJS exists but slow & Direct JS FFI; ideal for web \& Node deployment \\
Compile Times \& Tooling & Slower (especially with many extensions); Stack/Cabal & Faster; Spago/Purs tooling \\
Learning Curve \& Teams & Steep due to extensions and laziness & Gentler; closer to JavaScript developer mindset \\
Production Examples & CollegeVine, financial services, backend platforms & Web applications, startups with JS-heavy stacks \\
Simplicity (key advantage) & Complex: many extensions, laziness, advanced monads & Simple: no extensions, explicit forall, single canonical path for most features \\
\hline
\end{tabular}

For BEAM-centric production (scalable real-time systems), an Elixir-based reimplementation dramatically outperforms both originals by marrying strong typing with Erlang's runtime strengths.

\newpage
\section{Hindley-Milner Internals}

\subsection{Hindley-Milner Foundations}

Both languages (and Phi) share core HM features:

\begin{itemize}
\item  Automatic inference
\item  Parametric polymorphism (\texttt{id :: a -> a})
\item  Unification \& generalization
\item  Principal types
\item  Type classes for ad-hoc polymorphism
\item  Monadic effects for purity
\end{itemize}

\subsection{Key Differences in HM Implementation}

\begin{tabular}{|l|p{4.6cm}|p{4.6cm}|}
\hline
Feature & Haskell & PureScript \\
\hline
Polymorphism Syntax & Implicit \texttt{forall} & Explicit \texttt{forall a.} (required) \\
Inference Scope & Across modules; polymorphic recursion (with extensions) & Similar, but stricter module rules (no orphan instances) \\
HM Extensions & Many: higher-kinded types (HKT), GADTs, type families, DataKinds & Built-in only: row polymorphism for extensible records (\{ | r \}); no native HKT \\
Type Classes \& Instances & Unnamed instances; orphan instances allowed & Named instances; orphan instances forbidden \\
Deriving & Automatic \texttt{deriving (Eq)} inside data declarations & Explicit standalone \texttt{derive instance} \\
Evaluation Impact & Laziness can hide inference problems until runtime & Strict evaluation → predictable unification behavior \\
Records & Syntactic sugar; label-based; row polymorphism via extensions & Native row polymorphism → structural typing, close to JS objects \\
\hline
\end{tabular}

Phi keeps the HM core clean and explicit like PureScript, while eliminating JS limitations and adding BEAM-native constructs (e.g., direct process spawning, pattern-matched messaging).

\subsection{Type System Conclusion}

PureScript provides the cleanest, most production-friendly HM variant among the three originals — with native row polymorphism but without higher-kinded types (HKT), keeping complexity low. Haskell is far more expressive (including full HKT support) but often over-complex. A full Elixir reimplementation (Phi) delivers the best of both worlds: minimal HM + BEAM's unmatched concurrency/distribution.

\newpage
\section{Recommendation}

If the goal is **BEAM interoperability** and building scalable, fault-tolerant, distributed systems (5G, IoT, edge computing, telecom-grade applications), neither pure PureScript nor Haskell is sufficient on its own.

\begin{itemize}
\item PureScript excels in simplicity and web/full-stack but lacks native BEAM runtime.
\item Haskell offers power but brings extension complexity and no first-class BEAM support.
\end{itemize}

The superior path is a **full reimplementation in Elixir** — as demonstrated by the active Phi project (https://github.com/synrc/phi). Phi reimplements PureScript-like syntax, HM type inference/checking, ADTs, monads, and row polymorphism directly in Elixir, compiling to BEAM bytecode.

This approach yields:

\begin{itemize}
\item Static type safety at compile time
\item Pure functional style without extension soup
\item Native access to Erlang processes, supervision, distribution, and hot upgrades
\item Leverage of the mature Elixir/Erlang ecosystem
\end{itemize}

Both PureScript and Haskell remain excellent for their niches, but for BEAM-first production, Phi-style Elixir reimplementation is the clear way forward — dramatically safer, more scalable, and more maintainable than classical imperative alternatives on the Erlang VM.

\end{document}
